\chapter{Cadrage du projet et cahier des charges}
\label{chap:cadrage}

Une étude de gravité demande de préciser la population concernée avant de choisir une
méthode. Un nombre de personnes blessées, une proportion parmi les usagers impliqués
et un risque par kilomètre parcouru répondent à des questions différentes. Ce chapitre
fixe les objectifs, les conventions et les résultats attendus du projet.

\section{Finalité et questions de recherche}
Le projet vise à construire une étude reproductible sur des données publiques, depuis
leur audit jusqu'à l'interprétation des résultats. Le socle porte sur la qualité des
données et les analyses statistiques. Le volet prédictif complète l'étude
par la comparaison d'une référence constante et d'une régression logistique sur une
année ultérieure.

Trois questions structurent la première partie. Comment les catégories de gravité se
répartissent-elles dans la population enregistrée~? Quelles différences observe-t-on
selon l'âge, la catégorie d'usager, la lumière et l'agglomération~? Quelle variabilité
ces différences présentent-elles sous un schéma de rééchantillonnage qui respecte les
usagers d'un même accident~? Les comparaisons retenues sont en nombre limité et leur
sélection après exploration est déclarée.

Le volet Machine Learning vise à mesurer un gain par rapport à une probabilité
constante, puis à examiner les erreurs et la qualité des probabilités estimées. Son
intérêt est méthodologique~: aucune nécessité opérationnelle d'un outil de scoring
n'a été démontrée à ce stade. Une future utilisation demanderait de définir une décision,
un utilisateur et l'instant auquel les variables seraient disponibles.

\section{Périmètre et unité d'analyse}
La fenêtre choisie couvre six millésimes consécutifs, de 2019 à 2024. Le périmètre
géographique est la France métropolitaine, Corse comprise. Les sources complètes sont
conservées et auditées avant la sélection de la métropole. Ce choix rend explicite la
population étudiée~; il ne suppose pas une homogénéité des conditions de circulation
entre tous les territoires retenus.

Une observation représente \textbf{un usager source impliqué dans un accident}.
Deux occupants d'un véhicule restent deux observations, même si leurs caractéristiques
sont identiques. L'identifiant de ligne créé pendant la préparation n'est pas une clé
de personne permettant de la suivre d'un accident ou d'une année à l'autre.

Les personnes indemnes de BAAC sont impliquées dans des accidents corporels enregistrés.
Elles ne forment pas un groupe témoin de trajets sans accident. Les probabilités étudiées
restent donc conditionnelles à la présence dans ces enregistrements. Une comparaison
du risque par déplacement demanderait des données d'exposition supplémentaires.

\section{Définition de la gravité étudiée}
La variable \code{grav} comporte quatre modalités dont le sens est donné par la
documentation ONISR~\cite{doc-onisr}. Le décès est défini sur le coup ou dans les
30 jours suivant l'accident~; la catégorie hospitalisée correspond à plus de 24 heures
d'hospitalisation. Le tableau~\ref{tab:cible} présente le regroupement retenu.

\begin{table}[htbp]\centering\small
\caption{Modalités de gravité et cible principale.}\label{tab:cible}
\begin{tabularx}{\textwidth}{@{}lLl@{}}\toprule
Code & Modalité originale & Cible $Y$\\\midrule
1 & Indemne & 0\\
2 & Tué & 1\\
3 & Blessé hospitalisé & 1\\
4 & Blessé léger & 0\\
Autre / inconnu & Information non exploitable & Manquante\\\bottomrule
\end{tabularx}
\sourcecalcul{documentation ONISR et convention du projet}
\end{table}

La cible est intitulée \textit{décès ou hospitalisation}. Ce choix exprime une issue
compréhensible et fournit davantage de cas positifs que les décès seuls pour examiner
les sous-populations. Il ne résulte pas d'une recherche du meilleur score prédictif.
Les quatre modalités originales restent disponibles et les décès seuls font l'objet
d'une analyse complémentaire.

Ce regroupement n'est pas une définition médicale des blessures graves. La documentation
signale une modification de la saisie de l'hospitalisation depuis 2018 et l'absence de
labellisation de cet indicateur depuis 2019~\cite{doc-onisr}. Cette limite concerne toute
interprétation de la cible, y compris lorsque les calculs sont techniquement corrects.
Une gravité inconnue ne reçoit jamais automatiquement la valeur zéro.

\section{Exigences de préparation et de qualité}
Une transformation doit être accompagnée d'une raison et d'un contrôle. La préparation
vise à préserver les observations retenues, à rendre les informations inexploitables
visibles et à éviter les corrections fondées sur une supposition.

\begin{table}[htbp]\centering\small\renewcommand{\arraystretch}{1.2}
\caption{Exigences du cahier des charges pour la première partie.}\label{tab:exigences}
\begin{tabularx}{\textwidth}{@{}LL@{}}\toprule
Exigence & Critère de vérification\\\midrule
Traçabilité & Fichier source, rang de l'enregistrement, manifeste et empreintes SHA-256.\\
Clés cohérentes & Identifiants préparés en texte, clés complètes et absence de collisions.\\
Conservation des usagers & Identifiants et effectifs identiques avant et après les jointures.\\
Valeurs inconnues explicites & Règles par variable, sans remplacement global des zéros.\\
Lieux multiples maîtrisés & Une ligne de résumé par accident et désaccords signalés.\\
Analyses vérifiables & Effectifs, dénominateurs, proportions et origine des figures précisés.\\
Reproductibilité & Lecture des exports conforme aux types et valeurs de la table préparée.\\\bottomrule
\end{tabularx}
\end{table}

Les livrables sont les sept notebooks, le module ML partagé, les deux tables préparées,
les bilans de contrôle, les règles de lecture, les tableaux et figures, ainsi que
le présent rapport avec ses sources LaTeX. Le README, le fichier de dépendances et
le Makefile décrivent et facilitent l'installation. Une base SQL, une application web ou une extension en recherche
opérationnelle ne conditionnent pas la réussite de cette première version.

\section{Organisation temporelle et état du travail}
Les analyses destinées à guider les modèles sont limitées à 2019-2022. La validation
prédictive utilise 2023 et l'évaluation finale 2024. Cette dernière année a déjà été
examinée dans la découverte et les audits~: le test est donc chronologique et
rétrospectif, sans statut de jeu resté entièrement inconnu. Ses performances ne servent
pas à choisir les variables ou à régler les modèles.

Les règles de préparation sont déterministes. Aucune imputation par une moyenne des
six années, standardisation apprise ou sélection fondée sur les performances n'a été
effectuée. L'encodage du volet ML est ajusté sur l'apprentissage uniquement.

La première partie réalise les objectifs de préparation, d'exploration et de
comparaison statistique. La seconde examine les caractéristiques simultanément et
compare les probabilités estimées à la référence. Aucun seuil arbitraire de performance
n'est imposé~: les gains comme les limites font partie des résultats à analyser.

\FloatBarrier
\bilan{Le cahier des charges fixe une population, une unité et une cible explicites.
La réussite repose d'abord sur des données traçables et des conclusions correctement
délimitées. Le volet prédictif évalue ensuite ce que les informations disponibles
permettent effectivement d'estimer.}
